\QHeading{第1题　口算。}
\par\noindent


\begin{oralgrid}
\OralItem{1}{3.8+0.5}{4.3}
\OralItem{2}{3.6-1.4}{2.2}
\OralItem{3}{4.6+1.4}{6}
\OralItem{4}{3.7+2.1}{5.8}
\OralItem{5}{2.5-1.5}{1}
\OralItem{6}{1.9+1.1}{3}
\OralItem{7}{5.2+3.8}{9}
\OralItem{8}{4.7-3.7}{1}
\OralItem{9}{4.2+5.8}{10}
\OralItem{10}{0.2+7.9}{8.1}
\OralItem{11}{7.3-0.9}{6.4}
\OralItem{12}{20.9-9.7}{11.2}\end{oralgrid}
\par

\QHeading{口算思路补充（对应上表题号）}

\Answer{上面列出得数，下面按小题给出不同计算思路。计数单位法适用于小数，整数题可比较凑整、拆分或商不变等方法。}
\MethodTitle{第1小题：$3.8+0.5$}

\MethodBlock{口算思路一：按相同计数单位计算}{先看成若干个十分之一，按整数加减，最后还原成小数。}{\begin{aligned}[t]&\phantom{={}}3.8+0.5\\
&=(38+5)\div 10\\
&=43\div 10\\
&=4.3\end{aligned}}
\MethodTitle{第2小题：$3.6-1.4$}

\MethodBlock{口算思路一：按相同计数单位计算}{先看成若干个十分之一，按整数加减，最后还原成小数。}{\begin{aligned}[t]&\phantom{={}}3.6-1.4\\
&=(36-14)\div 10\\
&=22\div 10\\
&=2.2\end{aligned}}

\MethodBlock{口算思路二：补差法}{从减数先补到一个整数，再补到被减数，两段差相加。}{\begin{aligned}[t]&\phantom{={}}3.6-1.4\\
&=(2-1.4)+(3.6-2)\\
&=0.6+1.6\\
&=2.2\end{aligned}}
\MethodTitle{第3小题：$4.6+1.4$}

\MethodBlock{口算思路一：按相同计数单位计算}{先看成若干个十分之一，按整数加减，最后还原成小数。}{\begin{aligned}[t]&\phantom{={}}4.6+1.4\\
&=(46+14)\div 10\\
&=60\div 10\\
&=6\end{aligned}}
\MethodTitle{第4小题：$3.7+2.1$}

\MethodBlock{口算思路一：按相同计数单位计算}{先看成若干个十分之一，按整数加减，最后还原成小数。}{\begin{aligned}[t]&\phantom{={}}3.7+2.1\\
&=(37+21)\div 10\\
&=58\div 10\\
&=5.8\end{aligned}}
\MethodTitle{第5小题：$2.5-1.5$}

\MethodBlock{口算思路一：按相同计数单位计算}{先看成若干个十分之一，按整数加减，最后还原成小数。}{\begin{aligned}[t]&\phantom{={}}2.5-1.5\\
&=(25-15)\div 10\\
&=10\div 10\\
&=1\end{aligned}}

\MethodBlock{口算思路二：补差法}{从减数先补到一个整数，再补到被减数，两段差相加。}{\begin{aligned}[t]&\phantom{={}}2.5-1.5\\
&=(2-1.5)+(2.5-2)\\
&=0.5+0.5\\
&=1\end{aligned}}
\MethodTitle{第6小题：$1.9+1.1$}

\MethodBlock{口算思路一：按相同计数单位计算}{先看成若干个十分之一，按整数加减，最后还原成小数。}{\begin{aligned}[t]&\phantom{={}}1.9+1.1\\
&=(19+11)\div 10\\
&=30\div 10\\
&=3\end{aligned}}
\MethodTitle{第7小题：$5.2+3.8$}

\MethodBlock{口算思路一：按相同计数单位计算}{先看成若干个十分之一，按整数加减，最后还原成小数。}{\begin{aligned}[t]&\phantom{={}}5.2+3.8\\
&=(52+38)\div 10\\
&=90\div 10\\
&=9\end{aligned}}

\MethodBlock{口算思路二：凑整再减去多加的数}{把第二个数凑成整数，再减去多加的部分。}{\begin{aligned}[t]&\phantom{={}}5.2+3.8\\
&=5.2+4-0.2\\
&=9.2-0.2\\
&=9\end{aligned}}
\MethodTitle{第8小题：$4.7-3.7$}

\MethodBlock{口算思路一：按相同计数单位计算}{先看成若干个十分之一，按整数加减，最后还原成小数。}{\begin{aligned}[t]&\phantom{={}}4.7-3.7\\
&=(47-37)\div 10\\
&=10\div 10\\
&=1\end{aligned}}

\MethodBlock{口算思路二：补差法}{从减数先补到一个整数，再补到被减数，两段差相加。}{\begin{aligned}[t]&\phantom{={}}4.7-3.7\\
&=(4-3.7)+(4.7-4)\\
&=0.3+0.7\\
&=1\end{aligned}}
\MethodTitle{第9小题：$4.2+5.8$}

\MethodBlock{口算思路一：按相同计数单位计算}{先看成若干个十分之一，按整数加减，最后还原成小数。}{\begin{aligned}[t]&\phantom{={}}4.2+5.8\\
&=(42+58)\div 10\\
&=100\div 10\\
&=10\end{aligned}}

\MethodBlock{口算思路二：凑整再减去多加的数}{把第二个数凑成整数，再减去多加的部分。}{\begin{aligned}[t]&\phantom{={}}4.2+5.8\\
&=4.2+6-0.2\\
&=10.2-0.2\\
&=10\end{aligned}}
\MethodTitle{第10小题：$0.2+7.9$}

\MethodBlock{口算思路一：按相同计数单位计算}{先看成若干个十分之一，按整数加减，最后还原成小数。}{\begin{aligned}[t]&\phantom{={}}0.2+7.9\\
&=(2+79)\div 10\\
&=81\div 10\\
&=8.1\end{aligned}}

\MethodBlock{口算思路二：凑整再减去多加的数}{把第二个数凑成整数，再减去多加的部分。}{\begin{aligned}[t]&\phantom{={}}0.2+7.9\\
&=0.2+8-0.1\\
&=8.2-0.1\\
&=8.1\end{aligned}}
\MethodTitle{第11小题：$7.3-0.9$}

\MethodBlock{口算思路一：按相同计数单位计算}{先看成若干个十分之一，按整数加减，最后还原成小数。}{\begin{aligned}[t]&\phantom{={}}7.3-0.9\\
&=(73-9)\div 10\\
&=64\div 10\\
&=6.4\end{aligned}}

\MethodBlock{口算思路二：补差法}{从减数先补到一个整数，再补到被减数，两段差相加。}{\begin{aligned}[t]&\phantom{={}}7.3-0.9\\
&=(1-0.9)+(7.3-1)\\
&=0.1+6.3\\
&=6.4\end{aligned}}
\MethodTitle{第12小题：$20.9-9.7$}

\MethodBlock{口算思路一：按相同计数单位计算}{先看成若干个十分之一，按整数加减，最后还原成小数。}{\begin{aligned}[t]&\phantom{={}}20.9-9.7\\
&=(209-97)\div 10\\
&=112\div 10\\
&=11.2\end{aligned}}

\MethodBlock{口算思路二：补差法}{从减数先补到一个整数，再补到被减数，两段差相加。}{\begin{aligned}[t]&\phantom{={}}20.9-9.7\\
&=(10-9.7)+(20.9-10)\\
&=0.3+10.9\\
&=11.2\end{aligned}}


\par\noindent

\QHeading{第2题　竖式计算。}

\par\noindent\makebox[\linewidth][c]{
\begin{tikzpicture}[x=1bp,y=-1bp,color=black,draw=black,text=black]
\begin{scope}\clip (61.0,249.0) rectangle (535.0,430.0);
\hdhead{60}{小数加减法——信息窗1}
\hdsection{1}{140}{口算。}
\q{103}{185}{180}{13}{$3.8+0.5=$}
\q{103}{230}{180}{13}{$2.5-1.5=$}
\q{103}{275}{180}{13}{$4.2+5.8=$}
\q{293}{185}{180}{13}{$3.6-1.4=$}
\q{293}{230}{180}{13}{$1.9+1.1=$}
\q{293}{275}{180}{13}{$0.2+7.9=$}
\q{482}{184}{180}{13}{$4.6+1.4=$}
\q{482}{229}{180}{13}{$5.2+3.8=$}
\q{482}{274}{180}{13}{$7.3-0.9=$}
\q{675}{179}{180}{13}{$3.7+2.1=$}
\q{675}{224}{180}{13}{$4.7-3.7=$}
\q{675}{269}{180}{13}{$20.9-9.7=$}
\hdsection{2}{332}{竖式计算。}
\q{139}{388}{20}{14}{$\vphantom{0}1$}
\q{159}{388}{20}{14}{$\vphantom{0}.$}
\q{178}{388}{20}{14}{$\vphantom{0}7$}
\q{198}{388}{20}{14}{$\vphantom{0}9$}
\q{139}{428}{20}{14}{$\vphantom{0}0$}
\q{159}{428}{20}{14}{$\vphantom{0}.$}
\q{178}{428}{20}{14}{$\vphantom{0}8$}
\q{198}{428}{20}{14}{$\vphantom{0}9$}
\q{110}{428}{28}{14}{$\vphantom{0}-$}
\begin{scope}[scale=.6666667,draw=black,text=black]\draw[draw=black,line width=.6pt] (102.000,464.000)--(218.000,464.000);\end{scope}
\q{327}{387}{20}{14}{$\vphantom{0}0$}
\q{347}{387}{20}{14}{$\vphantom{0}.$}
\q{367}{387}{20}{14}{$\vphantom{0}3$}
\q{387}{387}{20}{14}{$\vphantom{0}7$}
\q{327}{427}{20}{14}{$\vphantom{0}0$}
\q{347}{427}{20}{14}{$\vphantom{0}.$}
\q{367}{427}{20}{14}{$\vphantom{0}6$}
\q{387}{427}{20}{14}{$\vphantom{0}2$}
\q{298}{427}{28}{14}{$\vphantom{0}+$}
\begin{scope}[scale=.6666667,draw=black,text=black]\draw[draw=black,line width=.6pt] (288.000,464.000)--(409.000,464.000);\end{scope}
\q{519}{385}{20}{14}{$\vphantom{0}4$}
\q{539}{385}{20}{14}{$\vphantom{0}.$}
\q{559}{385}{20}{14}{$\vphantom{0}5$}
\q{579}{385}{20}{14}{$\vphantom{0}6$}
\q{519}{425}{20}{14}{$\vphantom{0}2$}
\q{539}{425}{20}{14}{$\vphantom{0}.$}
\q{559}{425}{20}{14}{$\vphantom{0}4$}
\q{579}{425}{20}{14}{$\vphantom{0}3$}
\q{489}{425}{28}{14}{$\vphantom{0}-$}
\begin{scope}[scale=.6666667,draw=black,text=black]\draw[draw=black,line width=.6pt] (480.000,462.000)--(602.000,462.000);\end{scope}
\q{711}{383}{20}{14}{$\vphantom{0}0$}
\q{731}{383}{20}{14}{$\vphantom{0}.$}
\q{752}{383}{20}{14}{$\vphantom{0}7$}
\q{774}{383}{20}{14}{$\vphantom{0}7$}
\q{711}{423}{20}{14}{$\vphantom{0}5$}
\q{731}{423}{20}{14}{$\vphantom{0}.$}
\q{752}{423}{20}{14}{$\vphantom{0}2$}
\q{684}{423}{28}{14}{$\vphantom{0}+$}
\begin{scope}[scale=.6666667,draw=black,text=black]\draw[draw=black,line width=.6pt] (675.000,460.000)--(795.000,460.000);\end{scope}
\q{134}{542}{20}{14}{$\vphantom{0}1$}
\q{154}{542}{20}{14}{$\vphantom{0}2$}
\q{176}{542}{20}{14}{$\vphantom{0}.$}
\q{197}{542}{20}{14}{$\vphantom{0}2$}
\q{218}{542}{20}{14}{$\vphantom{0}1$}
\q{154}{582}{20}{14}{$\vphantom{0}3$}
\q{176}{582}{20}{14}{$\vphantom{0}.$}
\q{197}{582}{20}{14}{$\vphantom{0}4$}
\q{218}{582}{20}{14}{$\vphantom{0}7$}
\q{106}{582}{28}{14}{$\vphantom{0}+$}
\begin{scope}[scale=.6666667,draw=black,text=black]\draw[draw=black,line width=.6pt] (99.000,618.000)--(236.000,618.000);\end{scope}
\q{326}{541}{20}{14}{$\vphantom{0}2$}
\q{347}{541}{20}{14}{$\vphantom{0}9$}
\q{368}{541}{20}{14}{$\vphantom{0}.$}
\q{388}{541}{20}{14}{$\vphantom{0}7$}
\q{326}{581}{20}{14}{$\vphantom{0}1$}
\q{347}{581}{20}{14}{$\vphantom{0}8$}
\q{368}{581}{20}{14}{$\vphantom{0}.$}
\q{388}{581}{20}{14}{$\vphantom{0}8$}
\q{297}{581}{28}{14}{$\vphantom{0}-$}
\begin{scope}[scale=.6666667,draw=black,text=black]\draw[draw=black,line width=.6pt] (287.000,617.000)--(409.000,617.000);\end{scope}
\q{516}{540}{20}{14}{$\vphantom{0}4$}
\q{537}{540}{20}{14}{$\vphantom{0}.$}
\q{558}{540}{20}{14}{$\vphantom{0}3$}
\q{516}{580}{20}{14}{$\vphantom{0}1$}
\q{537}{580}{20}{14}{$\vphantom{0}.$}
\q{558}{580}{20}{14}{$\vphantom{0}2$}
\q{579}{580}{20}{14}{$\vphantom{0}1$}
\q{488}{580}{28}{14}{$\vphantom{0}-$}
\begin{scope}[scale=.6666667,draw=black,text=black]\draw[draw=black,line width=.6pt] (479.000,616.000)--(602.000,616.000);\end{scope}
\q{711}{539}{20}{14}{$\vphantom{0}4$}
\q{731}{539}{20}{14}{$\vphantom{0}.$}
\q{752}{539}{20}{14}{$\vphantom{0}4$}
\q{774}{539}{20}{14}{$\vphantom{0}4$}
\q{711}{579}{20}{14}{$\vphantom{0}1$}
\q{731}{579}{20}{14}{$\vphantom{0}.$}
\q{752}{579}{20}{14}{$\vphantom{0}5$}
\q{774}{579}{20}{14}{$\vphantom{0}4$}
\q{684}{579}{28}{14}{$\vphantom{0}+$}
\begin{scope}[scale=.6666667,draw=black,text=black]\draw[draw=black,line width=.6pt] (675.000,616.000)--(795.000,616.000);\end{scope}
\hdsection{3}{701}{解决问题。}
\q{99}{751}{741}{12}{① 一个文具盒8.45元，一块橡皮1.25元。买一个文具盒和一块橡皮一}
\q{99}{793}{741}{12}{　共需要多少钱?}
\q{99}{958}{741}{12}{② 超市里每袋牛奶2.4元，每袋方便面1.8元，一瓶矿泉水比一袋方便}
\q{99}{1000}{741}{12}{　面便宜0.65元。一瓶矿泉水多少钱?}

\begin{scope}[color=answerred,draw=answerred,text=answerred]\ans{140.67}{122.67}{53.33}{10}{$\vphantom{0}4.3$}
\ans{140.67}{152.67}{53.33}{10}{$\vphantom{0}1$}
\ans{140.67}{182.67}{53.33}{10}{$\vphantom{0}10$}
\ans{266.00}{122.67}{53.33}{10}{$\vphantom{0}2.2$}
\ans{266.00}{152.67}{53.33}{10}{$\vphantom{0}3$}
\ans{266.00}{182.67}{53.33}{10}{$\vphantom{0}8.1$}
\ans{394.00}{122.67}{53.33}{10}{$\vphantom{0}6$}
\ans{394.00}{152.67}{53.33}{10}{$\vphantom{0}9$}
\ans{394.00}{182.67}{53.33}{10}{$\vphantom{0}6.4$}
\ans{523.33}{118.00}{53.33}{10}{$\vphantom{0}5.8$}
\ans{523.33}{148.00}{53.33}{10}{$\vphantom{0}1$}
\ans{532.00}{178.00}{53.33}{10}{$\vphantom{0}11.2$}
\ans{92.67}{314.00}{14.67}{11}{$\vphantom{0}0$}
\ans{106.00}{314.00}{14.67}{11}{$\vphantom{0}.$}
\ans{118.67}{314.00}{14.67}{11}{$\vphantom{0}9$}
\ans{132.00}{314.00}{14.67}{11}{$\vphantom{0}0$}
\ans{218.00}{313.33}{14.67}{11}{$\vphantom{0}0$}
\ans{231.33}{313.33}{14.67}{11}{$\vphantom{0}.$}
\ans{244.67}{313.33}{14.67}{11}{$\vphantom{0}9$}
\ans{258.00}{313.33}{14.67}{11}{$\vphantom{0}9$}
\ans{346.00}{312.00}{14.67}{11}{$\vphantom{0}2$}
\ans{359.33}{312.00}{14.67}{11}{$\vphantom{0}.$}
\ans{372.67}{312.00}{14.67}{11}{$\vphantom{0}1$}
\ans{386.00}{312.00}{14.67}{11}{$\vphantom{0}3$}
\ans{474.00}{312.00}{14.67}{11}{$\vphantom{0}5$}
\ans{487.33}{312.00}{14.67}{11}{$\vphantom{0}.$}
\ans{501.33}{312.00}{14.67}{11}{$\vphantom{0}9$}
\ans{516.00}{312.00}{14.67}{11}{$\vphantom{0}7$}
\ans{89.33}{416.67}{14.67}{11}{$\vphantom{0}1$}
\ans{102.67}{416.67}{14.67}{11}{$\vphantom{0}5$}
\ans{117.33}{416.67}{14.67}{11}{$\vphantom{0}.$}
\ans{131.33}{416.67}{14.67}{11}{$\vphantom{0}6$}
\ans{145.33}{416.67}{14.67}{11}{$\vphantom{0}8$}
\ans{217.33}{416.67}{14.67}{11}{$\vphantom{0}1$}
\ans{231.33}{416.67}{14.67}{11}{$\vphantom{0}0$}
\ans{245.33}{416.67}{14.67}{11}{$\vphantom{0}.$}
\ans{258.67}{416.67}{14.67}{11}{$\vphantom{0}9$}
\ans{344.00}{415.33}{14.67}{11}{$\vphantom{0}3$}
\ans{358.00}{415.33}{14.67}{11}{$\vphantom{0}.$}
\ans{372.00}{415.33}{14.67}{11}{$\vphantom{0}0$}
\ans{386.00}{415.33}{14.67}{11}{$\vphantom{0}9$}
\ans{474.00}{415.33}{14.67}{11}{$\vphantom{0}5$}
\ans{487.33}{415.33}{14.67}{11}{$\vphantom{0}.$}
\ans{501.33}{415.33}{14.67}{11}{$\vphantom{0}9$}
\ans{516.00}{415.33}{14.67}{11}{$\vphantom{0}8$}
\ans{66.00}{556.00}{472.00}{10}{$8.45+1.25=9.7$（元）。\\答：一共需要9.7元。}
\ans{66.00}{693.33}{258.00}{10}{$\setlength{\jot}{-2pt}\begin{aligned}[t]&\phantom{={}}1.8-0.65\\&=1.80-0.65\\&=1.15\text{（元）}\end{aligned}$}
\ans{340.00}{708.00}{202.67}{10}{答：一瓶矿泉水1.15元。}
\ans{516.00}{282.00}{14.67}{11}{$\vphantom{0}0$}
\ans{386.00}{360.00}{14.67}{11}{$\vphantom{0}0$}

\end{scope}
\end{scope}\pgfresetboundingbox\path[use as bounding box] (61.0,249.0) rectangle (535.0,430.0);
\end{tikzpicture}}\par



\par






\par\noindent

\Answer{
\begingroup\color{answerred}
\QHeading{计算过程：2. 竖式计算}

小数点对齐，相同数位相加减；位数不足时在小数末尾补零。竖式上方的小字表示退位后的数。每题按从右向左的顺序逐列计算。

\noindent\setlength{\parindent}{0pt}
（1）$1.79-0.89=0.9$

\Formula{\begin{array}{r@{\,}r@{}c@{}r@{}r}&\overset{\scriptscriptstyle0}{1}&.&\overset{\scriptscriptstyle17}{7}&9\\-&0&.&8&9\\\hline &0&.&9&0\end{array}}
百分位：$9-9=0$。\par
\smallskip{}十分位7减8不够，向个位退1。\par
\smallskip{}个位退位后：$1-1=0$。\par
\smallskip{}十分位得到：$7+10=17$。\par
\smallskip{}十分位相减：$17-8=9$。\par
\smallskip{}个位相减：$0-0=0$。\par
\smallskip{}得0.90，化简为0.9。\par
%
\setlength{\parindent}{0pt}
（2）$0.37+0.62=0.99$

\Formula{\begin{array}{r}0.37\\+0.62\\\hline0.99\end{array}}
百分位：$7+2=9$。\par
\smallskip{}十分位：$3+6=9$。\par
\smallskip{}个位：$0+0=0$。\par
\smallskip{}得0.99。\par
\par\medskip

\noindent\setlength{\parindent}{0pt}
（3）$4.56-2.43=2.13$

\Formula{\begin{array}{r}4.56\\-2.43\\\hline2.13\end{array}}
百分位：$6-3=3$。\par
\smallskip{}十分位：$5-4=1$。\par
\smallskip{}个位：$4-2=2$。\par
\smallskip{}得2.13。\par
%
\setlength{\parindent}{0pt}
（4）$0.77+5.2=5.97$

\Formula{\begin{array}{r}0.77\\+5.20\\\hline5.97\end{array}}
把5.2写成5.20。\par
\smallskip{}百分位：$7+0=7$。\par
\smallskip{}十分位：$7+2=9$。\par
\smallskip{}个位：$0+5=5$。\par
\smallskip{}得5.97。\par
\par\medskip

\noindent\setlength{\parindent}{0pt}
（5）$12.21+3.47=15.68$

\Formula{\begin{array}{r}12.21\\+3.47\\\hline15.68\end{array}}
百分位：$1+7=8$。\par
\smallskip{}十分位：$2+4=6$。\par
\smallskip{}个位：$2+3=5$。\par
\smallskip{}十位：$1+0=1$。\par
\smallskip{}得15.68。\par
%
\setlength{\parindent}{0pt}
（6）$29.7-18.8=10.9$

\Formula{\begin{array}{r@{\,}r@{}r@{}c@{}r}&2&\overset{\scriptscriptstyle8}{9}&.&\overset{\scriptscriptstyle17}{7}\\-&1&8&.&8\\\hline&1&0&.&9\end{array}}
十分位7减8不够，向个位退1。\par
\smallskip{}个位退位后：$9-1=8$。\par
\smallskip{}十分位得到：$7+10=17$。\par
\smallskip{}十分位相减：$17-8=9$。\par
\smallskip{}个位相减：$8-8=0$。\par
\smallskip{}十位相减：$2-1=1$。\par
\smallskip{}得10.9。\par
\par\medskip

\noindent\setlength{\parindent}{0pt}
（7）$4.3-1.21=3.09$

\Formula{\begin{array}{r@{\,}r@{}c@{}r@{}r}&4&.&\overset{\scriptscriptstyle2}{3}&\overset{\scriptscriptstyle10}{0}\\-&1&.&2&1\\\hline&3&.&0&9\end{array}}
先把4.3写成4.30。\par
\smallskip{}百分位0减1不够，向十分位退1。\par
\smallskip{}十分位退位后：$3-1=2$。\par
\smallskip{}百分位得到：$0+10=10$。\par
\smallskip{}百分位相减：$10-1=9$。\par
\smallskip{}十分位相减：$2-2=0$。\par
\smallskip{}个位相减：$4-1=3$。\par
\smallskip{}得3.09。\par
%
\setlength{\parindent}{0pt}
（8）$4.44+1.54=5.98$

\Formula{\begin{array}{r}4.44\\+1.54\\\hline5.98\end{array}}
百分位：$4+4=8$。\par
\smallskip{}十分位：$4+5=9$。\par
\smallskip{}个位：$4+1=5$。\par
\smallskip{}得5.98。\par
\par\medskip


\endgroup

}

\QHeading{第3题　解决问题。}

\Problem{① 一个文具盒8.45元，一块橡皮1.25元。买一个文具盒和一块橡皮一共需要多少钱?}
\par


\MethodTitle{方法一}
\Answer{$8.45+1.25=9.7$（元）。\par 答：一共需要9.7元。}

\MethodTitle{方法二：先凑1元，再加余下的0.25元}
\Answer{把1.25元拆成1元和0.25元。}
\Answer{\Formula{\begin{aligned}[t]&\phantom{={}}8.45+1+0.25\\
&=9.45+0.25\\
&=9.7\text{（元）}\end{aligned}}}
\Answer{答：一共需要9.7元。}
\Problem{② 超市里每袋牛奶2.4元，每袋方便面1.8元，一瓶矿泉水比一袋方便面便宜0.65元。一瓶矿泉水多少钱?}

\MethodTitle{方法一}
\Answer{\Formula{\begin{aligned}[t]&\phantom{={}}1.8-0.65\\&=1.80-0.65\\&=1.15\text{（元）}\end{aligned}}}
\Answer{答：一瓶矿泉水1.15元。}

\MethodTitle{方法二：多减后加回}
\Answer{把0.65元先按0.7元减，多减了0.05元，再加回来。}
\Answer{\Formula{\begin{aligned}[t]&\phantom{={}}1.8-0.65\\
&=1.8-0.7+0.05\\
&=1.1+0.05\\
&=1.15\text{（元）}\end{aligned}}}
\Answer{答：一瓶矿泉水1.15元。}

